\documentclass[11pt, a4paper]{article}

% --- Paquets de base (Encodage et Langue) ---
\usepackage[utf8]{inputenc}
\usepackage[T1]{fontenc}
\usepackage[french]{babel}

% --- Design et Graphisme ---
\usepackage[default]{lato} % Police sans-serif moderne pour tout le document
\usepackage{xcolor}
\usepackage{graphicx}
\usepackage{geometry}
\usepackage{titlesec}
\usepackage{setspace}
\usepackage{hyperref}

% --- Couleurs et Marges ---
\definecolor{naturegreen}{RGB}{46, 125, 50} % Un beau vert "Écologie"
\definecolor{darkgray}{RGB}{50, 50, 50}

\geometry{
    top=3cm, 
    bottom=3cm, 
    left=2.5cm, 
    right=2.5cm,
    headheight=15pt
}

% --- Personnalisation des titres ---
\titleformat{\section}
  {\color{naturegreen}\Large\bfseries}{\thesection}{1em}{}[\titlerule]

\titleformat{\subsection}
  {\color{darkgray}\large\bfseries}{\thesubsection}{1em}{}

% --- En-tête et Pied de page ---
\usepackage{fancyhdr}
\pagestyle{fancy}
\fancyhf{}
\lhead{\textcolor{gray}{Rapport de Stage M1}}
\rhead{\textcolor{gray}{Écologie Tropicale}}
\cfoot{\thepage}
\renewcommand{\headrulewidth}{0.4pt}

% --- Variables du rapport ---
\newcommand{\monTitre}{Dynamique de la Biodiversité Ligneuse en Forêt Karstique}
\newcommand{\monNom}{Prénom NOM}
\newcommand{\monLabo}{Laboratoire d'Écologie des Systèmes Tropicaux}

% --- Début du document ---
\begin{document}

% --- Page de Garde Épurée ---
\begin{titlepage}
    \centering
    \vspace*{2cm}
    {\large \monLabo \par}
    \vspace{4cm}
    {\Huge \bfseries \color{naturegreen} \monTitre \par}
    \vspace{1cm}
    \rule{0.5\textwidth}{1pt}
    \vspace{2cm}
    
    {\Large \textbf{\monNom} \par}
    \vspace{1cm}
    {\large Sous la direction de \textbf{Dr. Dejun LI}}
    
    \vfill
    
    {\large Année Universitaire 2025 -- 2026 \par}
\end{titlepage}

\newpage
\onehalfspacing % Interligne agréable pour la lecture

% --- Résumé (Abstract) ---
\section*{Résumé}
\textit{
Ce rapport examine comment la diversité des espèces d'arbres influence la persistance du carbone organique du sol (SOC). Contrairement aux idées reçues, une diversité accrue peut diminuer le ratio entre le carbone associé aux minéraux (MAOC) et le carbone particulaire (POC) tout en augmentant la stabilité globale du pool de carbone. Cette étude démontre que la biodiversité forestière reconfigure activement les voies de stabilisation du carbone.
}

\vspace{2em}
\textbf{Mots-clés :} Écologie tropicale, Karst, Carbone organique du sol, Biodiversité.

\newpage
\tableofcontents
\newpage

\section{Introduction}
Le carbone organique du sol (SOC) constitue le plus grand réservoir de carbone terrestre. Comprendre les facteurs qui améliorent sa stabilité est crucial pour prédire les scénarios de changement climatique futurs. La diversité végétale est un levier majeur pour renforcer ces stocks.

\section{Matériels et Méthodes}
L'étude a été réalisée dans une forêt karstique subtropicale (Réserve de Mulun, Chine). Nous avons utilisé un gradient naturel de diversité allant de 4 à 61 espèces par parcelle.

\subsection{Analyses du sol}
Le sol a été échantillonné sur une profondeur de 0-10 cm. Les fractions de carbone ont été isolées par tamisage humide à 53 $\mu$m pour séparer le MAOC du POC.

\section{Résultats}
L'augmentation de la diversité des espèces d'arbres mène à une accumulation significative de SOC total. Cependant, elle induit une baisse marquée du ratio MAOC:POC. Cette reconfiguration est étroitement liée à la présence de carbonate de calcium ($CaCO_3$) qui protège physiquement le carbone particulaire.

\section{Conclusion}
La diversité des arbres ne se contente pas de stocker plus de carbone ; elle rend ce stockage plus résilient en activant des mécanismes de protection physique et microbien spécifiques.

\end{document}